% !TEX root = ../main.tex
% !TEX program = lualatex
\documentclass[../main.tex]{subfiles}
\IfSubfilesClassLoaded{\externaldocument{\subfix{../build/main}}}{}
% =============================================================================
% File: 21_Fleet_Maintenance_Policy_and_Implementation.tex
% =============================================================================
\begin{document}
\IfSubfilesClassLoaded{\chapter{EDO Study Guide}}{}
\section{Fleet Maintenance Policy and Implementation}

\subsection{Policy Baseline}
\begin{description}
  \item[\textbf{\ac{opnavinst} 4700.7M.}] Establishes the maintenance policy for United States Navy ships: maintain the fleet at the highest practical readiness, perform work at the lowest capable level, and apply \ac{rcm}/\ac{cbm} methods to manage Total Ownership Cost~\autocite{OPNAVINST4700-7M}.
\item[\textbf{\ac{opnavnote} 4700.}] Publishes \ac{cno} availability intervals, durations, man-day targets, and Repair Days for each platform, providing the authoritative baseline for the Long-Range Maintenance Schedule (\ac{lrms})~\autocite{OPNAVNOTE4700-24}.
\item[\textbf{\ac{jfmm} and \ac{cbmp}.}] \ac{usff}/\ac{pacflt}INST~4790.3 (\ac{jfmm}) assigns responsibilities to \ac{tycom}s, \ac{rmmco}s, and \ac{nsa}s, while \ac{opnavinst}~4790.16 issues \ac{cbmp} policy so sensors and analytics can replace purely time-based maintenance when justified~\autocite{COMFLTFORCOMINST4790-3,OPNAVINST4790-16G}.
\end{description}

\begin{longtblr}[
    caption = {Maintenance policy map at a glance},
    label   = {tab:maintenance_policy_map}
  ]{
    colspec = {@{} L L L L @{}}
  }
  \toprule
  {Document} & {Owner} & {Decisions Touched} & {Compliance Talking Points} \\
  \midrule
  OPNAVINST 4700.7M & CNO (N9) & Sets fleet maintenance policy baseline (lowest capable level, RCM/CBM, color-of-money discipline) & Use as the anchor for deferral justifications, craft risk acceptance to its hierarchy, and cite when defending public/private yard mix and maintenance levels \\
  OPNAVNOTE 4700 & CNO (N9) & Publishes class-by-class availability intervals, durations, manday targets, and repair days that drive the Long-Range Maintenance Schedule & Treat as the authoritative schedule baseline; deviations require TYCOM/CNO alignment and supporting maintenance metrics \\
  JFMM (COMFLTFORCOMINST 4790.3) & USFF/PACFLT & Defines execution standards for work control, assessments, certifications, and availability gate reviews & Align DFS/TSO closure, QA evidence, and RMMCO coordination to JFMM chapters; reference for watch team and AIT roles \\
  CBMP (OPNAVINST 4790.16) & CNO (N4) & Governs condition-based maintenance approvals, data pedigree, and sensor analytics used to replace time-directed tasks & Document the analytics package (sensor health, thresholds, reliability data) and secure CBMP authority approval before removing directed tasks \\
  \bottomrule
\end{longtblr}

The entire maintenance enterprise hinges on tying ship assessments to these references; citing a different standard during a board (e.g., \ac{oem} manual) without reconciling it to the \ac{jfmm} and policy hierarchy is a red flag.

\subsection{Strategy Frameworks}
\begin{description}
\item[\textbf{\ac{ofrp}.}] OPNAVINST~3000.15 drives Fleet Response Plans that align readiness, training, and maintenance phases; maintenance windows in the \ac{ofrp} billet plan are how \acp{tycom} defend availability lengths and sequencing~\autocite{OPNAVINST3000-15K}.
\item[\textbf{\ac{sfrm}.}] \ac{surflant}/\ac{surfpac} 3502-series guidance aligns training milestones (Basic, Advanced, Integrated) with maintenance health metrics so a ship cannot progress if critical maintenance or material conditions remain open~\autocite{CNSP-CNSLINST3502-7C}.
\item[\textbf{\ac{nsssy}.}] \ac{nsssy} injects commercial best practices (Theory of Constraints, value-stream management, digital dashboards) into the four public yards with the goal of 80\% on-time \acp{ssn} deliveries and no idle submarine backlog by FY26~\autocite{edo-4-1-3-fleet-maint-policy-2025}.
\end{description}

\subsection{Availability Approval and Execution Discipline}
\begin{enumerate}
\item \textbf{Baseline work.} \ac{surfmepp}/\ac{submepp} compile \ac{cmp} work, \acp{tycom} repairs, and modernization candidates into the \ac{bawp} at A-36 months; work then flows through the \ac{cno} scheduling conference for final approval~\autocite{OPNAVNOTE4700-24,edo-4-1-3-fleet-maint-policy-2025}.
\item \textbf{Directed vs.\ condition-based maintenance.} Directed Maintenance Tasks (\ac{dmt}) must be executed regardless of condition (e.g., structural checks mandated by \ac{subsafe}), while \ac{cbm} tasks rely on evidence such as oil analysis, smart sensors, or \ac{tsra} discoveries to justify deferral~\autocite{OPNAVINST4700-7M,OPNAVINST4790-16G}.
\item \textbf{Approval boards.} \acp{tycom} Availability Work Package Reviews validate that all \acp{dfs} and \acp{tso} have closure plans before endorsing the package for \ac{cno} approval~\autocite{COMFLTFORCOMINST4790-3}.
\item \textbf{Metrics.} Schedule adherence, manday performance, and quality escape metrics (e.g., repeat maintenance, Objective Quality Evidence (\ac{oqe}) rejects) appear in \ac{usff}/\ac{pacflt} readiness briefs and feed the next \ac{lrms} update~\autocite{edo-4-1-3-fleet-maint-policy-2025}.
\end{enumerate}

\subsection{Trends, Risks, and Talking Points}
\begin{itemize}
\item \textbf{Industrial base mix.} Navy policy seeks roughly a 50/50 public-private mix for surface depot work, with \acp{cvn} and \acp{ssn}/\acp{ssbn} remaining predominantly public-yard to retain nuclear competencies~\autocite{OPNAVINST4700-7M}.
\item \textbf{Manpower constraints.} \ac{nsssy} reviews highlight persistent shortages in key trades (nuclear welders, radiological controls, rigging) that drive docking delays; referencing \ac{nsssy} corrective actions shows awareness of current risk mitigations~\autocite{edo-4-1-3-fleet-maint-policy-2025}.
\item \textbf{Digital thread.} The \ac{nde} modernization module is now the authoritative source for installed configuration, allowing planners to reconcile modernization with maintenance data without resorting to ad-hoc spreadsheets~\autocite{TS9090-310}.
\item \textbf{Fiscal discipline.} Maintenance execution must honor color-of-money boundaries (\ac{ocmf} vs.\ \ac{opn}/\ac{scn}); unauthorized obligation adjustments are \ac{ada} violations~\autocite{DoDFMR-Vol3,FAR}.
\end{itemize}

\IfSubfilesClassLoaded{
  \RenewDocumentCommand{\entryname}{}{\textbf{\color{Modern} Acronym}}
  \RenewDocumentCommand{\descriptionname}{}{\textbf{\color{Modern} Definition}}
  \printnoidxglossary[
    type=\acronymtype,
    title=Acronyms,
    style=long-booktabs
  ]}{}
\end{document}
